% !TEX encoding = UTF-8 Unicode

\chapter{Grundlagen und verwandte Arbeiten}
\label{cha:related_work}

\blindmathtrue
\Blindtext

\begin{figure}[h!]
\centering
\subfloat[fig 1%
	\label{fig:figID_1}]{%
	\includegraphics[width=0.35\linewidth]{img/PDF_1}
}
\quad\quad\quad
\subfloat[fig 2%
	\label{fig:figID_2}]{%
	\includegraphics[width=0.35\linewidth]{img/PDF_2}
}
\\
\subfloat[fig 3%
	\label{fig:figID_3}]{%
	\includegraphics[width=0.35\linewidth]{img/PDF_3}
}
\quad\quad\quad
\subfloat[fig 4%
	\label{fig:figID_4}]{%
	\includegraphics[width=0.35\linewidth]{img/PDF_4}
}
\vspace{2mm}
\hrule
\caption{In der Bildunterschrift wird beschrieben, was in den einzelnen Bildern zu sehen ist. Hier werden vier PDF-Dateien geladen.}
\label{fig:figID}
\end{figure}


\Blindtext[2]


\begin{figure}[h!]
	\centering
	\includegraphics[width=0.4\linewidth]{img/EPS}
	\vspace{2mm}
	\hrule
	\caption{Auch EPS-Dateien können geladen werden.}
	\label{fig:fig_eps}
\end{figure}


\Blindtext

Zu diesem Thema gibt es viele Quellen \Cite{Report2019,Conference2019}  und zu diesem noch mehr wichtige Quellen \Cite{Article2019,Book2019}.

%\input{chp/theory_related_work/...}
%\input{chp/theory_related_work/...}
%\input{chp/theory_related_work/...}
